\chapter{概率母函数}\label{ch:b2:genfun}

\cref{ch:b2:powerseries}中的幂级数再次登场，
带着概率论的任务：对一个 $\N$ 值随机变量，
配以系数为 $\P(X = n)$ 的幂级数。这一
\emph{母函数}把独立随机变量之和
化为乘积，矩变为 $1$ 处的导数，而困难的
组合恒等式化为一行乘法。本章以两个
压轴例子收束全书：稀有事件的泊松近似，
以及分支过程的灭绝判据 ---
一个真正无穷的概率计算，完全由一条凸曲线的
几何来解决。

\section{定义与基本性质}

\begin{definition}[概率母函数]\label{def:b2:genfun:pgf}
设 $X$ 为 $\N$ 值随机变量，$p_n = \P(X = n)$。
称 $X$ 的\emph{概率母函数}\index{概率母函数}
为下列
幂级数之和：
\[
G_X(t) = \E\bigl(t^X\bigr) = \sum_{n=0}^{\infty} p_n\,t^n .
\]
\end{definition}

\begin{example}[基本手法]\label{ex:b2:genfun:firstreflexes}
常值随机变量 $X = c$ 的母函数为 $G_X(t) = t^c$；平移
满足 $G_{X+c}(t) = t^c\,G_X(t)$；在特殊点求值，
无需展开即可读出信息：$G_X(0) = \P(X
= 0)$，$G_X(1) = 1$，以及 $G_X(-1) = \P(X\text{ 为偶数}) -
\P(X\text{ 为奇数})$，这一宇称平衡被用于
\cref{exo:b2:genfun:10}。这些一行式结论
在下文中处处默认使用 --- 而求值 $G_X(0)$ 恰好
正是本章末尾从迭代母函数中提取灭绝概率的
途径。
\end{example}

\begin{proposition}[收敛半径与基本性质]\label{prop:b2:genfun:radius}
定义 $G_X$ 的级数，其收敛半径 $\geq 1$；
$G_X$ 在 $\intcc{-1}{1}$ 上有定义且连续，
为 $\mathcal{C}^\infty$ 函数（在开区间 $\intoo{-1}{1}$ 上），且 $G_X(1) = 1$，
在该区间上 $\abs{G_X(t)} \leq 1$。此外，$G_X$ 决定
了 $X$ 的分布：
\[
p_n = \frac{G_X^{(n)}(0)}{n!} .
\]
\end{proposition}

\begin{proof}
由 $\sum p_n = 1$ 收敛可知，项 $p_n\,1^n$ 有界，
故收敛半径 $\geq 1$（阿贝尔引理，
见\cref{ch:b2:powerseries}）；在 $t = \pm1$ 处级数绝对收敛
（$\sum p_n = 1$ 占优）；更进一步，在整个
区间 $\intcc{-1}1$ 上，
\[
\sup_{\abs t\leq1}\,\abs{p_nt^n} = p_n
\quad\text{其中}\quad \sum_np_n < \infty :
\]
级数在 $\intcc{-1}1$ 上\emph{正规收敛}，故
其和在该区间上连续，
（由\cref{thm:b2:funcseq:weierstrass,thm:b2:funcseq:continuity}可知）。区间内部的光滑性与
系数公式都属于幂级数的一般理论；系数
可以恢复，故母函数相同的两个随机变量
具有相同的分布。
\end{proof}

\begin{example}[经典分布]\label{ex:b2:genfun:classical}
\leavevmode
\begin{itemize}
\item 伯努利分布 $\mathcal{B}(p)$：$G(t) = 1 - p + pt$。
\item 二项分布 $\mathcal{B}(n, p)$：$G(t) = \sum_k \binom nk
(pt)^k(1-p)^{n-k} = (1 - p + pt)^n$（二项式定理）。
\item 几何分布 $\mathcal{G}(p)$：$G(t) =
\sum_{k\geq1}(1-p)^{k-1}p\,t^k = \dfrac{pt}{1 - (1-p)t}$（收敛半径
$\frac{1}{1-p} > 1$）。
\item 泊松分布 $\mathcal{P}(\lambda)$：$G(t) =
\sum_k e^{-\lambda}\frac{(\lambda t)^k}{k!} = e^{\lambda(t - 1)}$
（收敛半径 $\infty$）。
\end{itemize}
\end{example}

\begin{example}[母函数的
积分]\label{ex:b2:genfun:integraltrick}
$G_X$ 在 $1$ 处的导数给出正矩；而
\emph{积分}给出一个负矩。由 $\int_0^1t^k\dd t
= \frac1{k+1}$，并逐项积分（
$\intcc01$ 上的正规收敛）：
\[
\int_0^1G_X(t)\,\dd t = \sum_{k\geq0}\frac{\P(X =
k)}{k+1} = \E\Bigl(\frac1{1+X}\Bigr).
\]
当 $X \sim \mathcal P(\lambda)$ 时：
\[
\E\Bigl(\frac1{1+X}\Bigr) =
\int_0^1\eu^{\lambda(t-1)}\,\dd t = \frac{1 -
\eu^{-\lambda}}{\lambda},
\]
一行即可重新得出
\cref{ex:b2:randomvar:transferex}中的级数计算。母函数
是一件双向工具：在 $1$ 处求导以得
矩 $\E(X)$、$\E(X(X-1))$，在 $\intcc01$ 上积分，
以求 $\E\bigl(\frac1{1+X}\bigr)$ --- 同一个解析对象，
随问题所需，从任一方向加以考察。
\end{example}

\begin{example}[收敛半径恰为
1 的分布]\label{ex:b2:genfun:heavytail}
设 $\P(X = k) = \dfrac{6}{\pi^2k^2}$，$k \geq 1$ --- 这
是一个概率分布，依据巴塞尔恒等式
（见\cref{ex:b2:fourier:basel}）。其母函数
$G(t) = \frac6{\pi^2}\sum_{k\geq1}\frac{t^k}{k^2}$ 的收敛半径
恰为 $1$：一般的界“收敛半径 $\geq
1$”（见\cref{prop:b2:genfun:radius}）不能再改进。
其期望为
\[
\sum_{k\geq1}k\,\P(X = k) =
\frac6{\pi^2}\sum_{k\geq1}\frac1k = \infty :
\]
$G$ 在 $\intcc{-1}1$ 上连续，内部光滑，但其
导数在 $1^-$ 处发散 --- 图像
以竖直切线抵达点 $(1, 1)$。重尾
在母函数上可以\emph{几何地}看出，且只在
单点 $t = 1$ 处；下文的矩定理使这一对应
精确成立。
\end{example}

\begin{theorem}[由母函数
求矩]\label{thm:b2:genfun:moments}
$X$ 有期望，当且仅当 $G_X$ 在 $1^-$
处可微（左导数有限），且此时 $\E(X) = G_X'(1)$。
类似地，$X$ 有二阶矩，当且仅当 $G_X$
在 $1^-$ 处二阶可微，且此时
\[
\E\bigl(X(X - 1)\bigr) = G_X''(1),
\qquad
V(X) = G_X''(1) + G_X'(1) - G_X'(1)^2 .
\]
\end{theorem}

\begin{proof}
对 $t \in \intoo{0}{1}$，在收敛圆盘内逐项求导，
得 $G_X'(t) = \sum_{n\geq1} np_n t^{n-1}$，这是一个
系数非负的级数：$t \mapsto G_X'(t)$ 单调不减，
在 $\intoo{0}{1}$ 上，由部分和的单调收敛（或
非负系数情形的阿贝尔定理，
见\cref{ch:b2:powerseries}），得
\[
\lim_{t \to 1^-} G_X'(t)
= \sum_{n\geq1} n\,p_n \in \intcc{0}{+\infty} ,
\]
两边有限当且仅当另一边有限。有限时，由中值
定理夹逼差商 $\frac{G_X(1) -
G_X(t)}{1 - t}$ 介于 $G_X'$ 的值之间，于是 $G_X$
在 $1^-$ 处可微，且 $G_X'(1) = \sum np_n = \E(X)$（由
转移公式）。二阶情形的论断把上述论证
再推进一层：$G''_X(t) = \sum_{n\geq2}n(n-1)p_nt^{n-2}$ 在
$\intoo01$ 上单调不减，单调极限为
$\sum_nn(n-1)p_n = \E(X(X-1))$，当且仅当 $X$ 有
二阶矩时有限。方差公式由
柯尼希–惠更斯公式得：
\[
V(X) = \E(X^2) - \E(X)^2 = \E\bigl(X(X-1)\bigr) + \E(X) -
\E(X)^2 = G''_X(1) + G'_X(1) - G'_X(1)^2 .
\]
\end{proof}

\begin{example}
对泊松分布，$G'(t) = \lambda e^{\lambda(t-1)}$，故 $\E(X) = \lambda$；
$G''(1) = \lambda^2$，故 $V(X) = \lambda^2 + \lambda - \lambda^2 =
\lambda$ --- 即在\cref{ch:b2:randomvar}中的
计算，每个只需一行。
\end{example}

\begin{example}[泊松分布的众数]
\label{ex:b2:genfun:poissonmode}
$\P(X = k)$ 在何处最大？设 $X \sim \mathcal
P(\lambda)$，相邻权重可通过比值比较：
\[
\frac{\P(X = k+1)}{\P(X = k)} = \frac{\lambda}{k + 1} ,
\]
该比值大于 $1$，当 $k < \lambda - 1$ 时，
而小于 $1$（当 $k > \lambda - 1$ 时），故权重先升后降，
众数为 $\floor\lambda$（$\lambda - 1$ 与
$\lambda$ 并列为众数（$\lambda$ 为整数时）：对 $\lambda = 3$，
$\P(X = 2) = \P(X = 3) = \frac92\eu^{-3} \approx 0.224$）。
对系数作比值判别，往往是得到离散分布的
定性性质的最快途径，无需母函数——
但系数\emph{本身即是}母函数，
逐项读出即可。
\end{example}

\section{独立变量之和}

\begin{theorem}[乘性]\label{thm:b2:genfun:product}
设 $X$ 与 $Y$ 为独立的 $\N$ 值随机变量，则
\[
G_{X + Y}(t) = G_X(t)\,G_Y(t)
\qquad (\abs t \leq 1),
\]
由归纳法，$G_{X_1 + \dots + X_n} = \prod_i G_{X_i}$
对独立的 $X_1, \dots, X_n$ 成立。
\end{theorem}

\begin{proof}
两种证明都很有启发性。\emph{利用期望：}$t^X$ 与
$t^Y$ 是独立的有界变量，
由\cref{thm:b2:randomvar:product}可知
\[
G_{X+Y}(t) = \E\bigl(t^{X+Y}\bigr)
= \E\bigl(t^X t^Y\bigr)
= \E\bigl(t^X\bigr)\E\bigl(t^Y\bigr) .
\]
\emph{利用柯西乘积：}$X + Y$ 的分布即为卷积
$\P(X + Y = n) = \sum_{k=0}^n \P(X = k)\P(Y = n - k)$，
而对绝对收敛级数的柯西乘积定理
（见\cref{ch:b2:series}）使两个幂级数
沿此卷积精确相乘。
\end{proof}

\begin{example}[经典分布的稳定性]\label{ex:b2:genfun:stability}
具有相同 $p$ 的独立二项分布可以相加：$(1 - p + pt)^m(1 -
p + pt)^n = (1 - p + pt)^{m+n}$，故 $\mathcal{B}(m, p) +
\mathcal{B}(n, p) = \mathcal{B}(m + n, p)$ --- 特别地，$n$ 个独立
伯努利变量之和服从二项分布，重新证明了
成功次数的分布。独立的泊松变量可以相加：
$e^{\lambda(t-1)}e^{\mu(t-1)} = e^{(\lambda + \mu)(t-1)}$，故
$\mathcal{P}(\lambda) + \mathcal{P}(\mu) = \mathcal{P}(\lambda +
\mu)$ --- 正是
\cref{exo:b2:randomvar:2}中的卷积计算，此处无需再算。
\end{example}

\begin{example}[两颗骰子，一个多项式
的平方]\label{ex:b2:genfun:twodice}
对一颗均匀骰子，$G(t) = \frac{t + t^2 + \dots + t^6}{6}$；
对两颗骰子之和，
\[
G(t)^2 = \frac{1}{36}\bigl(t^2 + 2t^3 + 3t^4 + 4t^5 + 5t^6 +
6t^7 + 5t^8 + 4t^9 + 3t^{10} + 2t^{11} + t^{12}\bigr) :
\]
骰子点数和的三角形分布（$7$ 为众数，
概率 $\frac6{36} = \frac16$），由一个多项式的平方
直接读出，这种平方一生只需展开一次。
卷积公式需要十一个
独立的计数论证，而母函数
同时完成所有计数，因为多项式相乘
\emph{即是}系数的卷积。这种机械的转换——
分布化为系数，和化为积——
正是本章的全部核心思想，而
\cref{exo:b2:genfun:11}将其推向出人意料的
西彻曼骰子。
\end{example}

\begin{example}[三个骰子与
系数提取]\label{ex:b2:genfun:threedice}
设 $S$ 为三个均匀骰子的点数之和，则 $\P(S = 10)$ 为
$t^{10}$ 在 $\bigl(\frac{t + \dots +
t^6}6\bigr)^3$ 中的系数。分解因式，借助二项式与
几何级数展开：
\[
\Bigl(\frac{t(1 - t^6)}{6(1 - t)}\Bigr)^{\!3}
= \frac{t^3}{216}\,\bigl(1 - 3t^6 + 3t^{12} -
t^{18}\bigr)\sum_{j\geq0}\binom{j+2}2t^j .
\]
要得 $t^{10}$ 的系数，需从乘积中取出 $t^7$：
即 $j = 7$ 时取 $1$，$j = 1$ 时取
$-3t^6$：
\[
\P(S = 10) = \frac{1}{216}\Bigl(\binom92 -
3\binom32\Bigr) = \frac{36 - 9}{216} = \frac{27}{216} =
\frac18 .
\]
直接逐一枚举 $27$ 个三元组易出错，而
代数运算是机械的，可推广到任意多个骰子 ---
$(1 - t^6)^3$ 中显现的容斥原理正在
自动完成分情形的讨论。
\end{example}

\begin{example}[由母函数读出
分布]\label{ex:b2:genfun:identify}
哪一个分布的母函数为 $G(t) = \dfrac1{2 - t}$？展开为幂
级数：
\[
\frac{1}{2 - t} = \frac12\cdot\frac1{1 - t/2}
= \sum_{k\geq0}\frac{t^k}{2^{k+1}} :
\]
系数非负且和为 $G(1) = 1$，故这是一个真正的
分布，$\P(X = k) = 2^{-(k+1)}$，取值于 $\N$ ---
即从 $0$ 开始的几何分布。由唯一性
（见\cref{prop:b2:genfun:radius}），没有别的分布具有相同的
$G$。由母函数识别分布是一项
值得反复练习的技能，正是它揭示了临界分支
迭代 $G_n(t) = \frac{n - (n-1)t}{n+1 - nt}$ 在
周末问题中，实为一个
以存活为条件的几何分布。
\end{example}

\begin{remark}
稳定性只朝一个方向成立：独立的泊松随机变量之和
仍为泊松分布，而差则不然 --- $X - Y$ 取
负值，故其根本没有母函数，其分布
（即斯凯拉姆分布）超出本章的
工具范围。同样，$\mathcal B(m, p) + \mathcal B(n, p')$ 在
$p \neq p'$ 时，\emph{并非}二项分布：乘积 $(1 - p +
pt)^m(1 - p' + p't)^n$ 有两个不同的根位置，而
每个二项分布的概率母函数都只有一个重根。
由根的模式读出稳定性，只是一个小小的预览，
展示多项式所蕴含的结构之丰富。
\end{remark}

\begin{remark}[单位根滤波]
在 $-1$ 处求值可区分偶数与奇数；在所有
$m$ 次单位根处求值则可分离每一剩余类：令
$\omega = \eu^{2\iu\pi/m}$,
\[
\P(X \equiv r \bmod m)
= \frac1m\sum_{j=0}^{m-1}\omega^{-jr}\,G_X(\omega^j),
\]
由于对 $\omega^{j(k-r)}$ 按 $j$ 求平均，得 $1$ 当 $k
\equiv r$ 时，否则为 $0$。例：设 $S$ 为
两个均匀骰子的点数之和，则每个 $G(\omega^j) = \frac16\sum_{k=1}^6
\omega^{jk} = -\frac16$，当 $j \neq 0$ 时（七个七次
单位根之和为零），故
\[
\P(7 \mid S) = \frac17\Bigl(1 +
6\cdot\frac1{36}\Bigr) = \frac16 ,
\]
与\cref{ex:b2:genfun:twodice}中的计数相吻合 ---
且该方法可推广到直接计数难以处理的
问题。
\end{remark}

\begin{theorem}[随机和：母函数形式的瓦尔德
恒等式]\label{thm:b2:genfun:compound}
设 $(X_k)_{k\geq1}$ 为相互独立、取值于 $\N$ 的随机变量，它们
具有相同的分布与母函数 $G_X$，又设 $N$ 为
取值于 $\N$ 且与所有 $X_k$ 独立的随机变量，其
母函数为 $G_N$。则随机和 $S = X_1 + \dots + X_N$（约定
$S = 0$，当 $N = 0$ 时）的母函数为
\[
G_S = G_N \circ G_X .
\]
特别地，若 $N$ 与 $X_1$ 存在期望，则 $\E(S) =
\E(N)\,\E(X_1)$.
\end{theorem}

\begin{proof}
以 $N$ 为条件（全概率公式，
见\cref{thm:b2:proba:bayes}）：对 $\abs t \leq 1$，
\[
G_S(t) = \sum_{n=0}^\infty \P(N = n)\,
\E\bigl(t^{X_1 + \dots + X_n}\bigr)
= \sum_{n=0}^\infty \P(N = n)\,G_X(t)^n
= G_N\bigl(G_X(t)\bigr),
\]
对每个固定的 $n$ 利用乘性，以及
整个二重族的可和性（$\abs{G_X(t)} \leq 1$）。求和次序的交换
即为可和族的富比尼定理（
\cref{ch:b2:series}）。在 $1^-$ 处求导，由链式法则及
\cref{thm:b2:genfun:moments}，得 $\E(S) = G_N'(G_X(1))\,G_X'(1)
= G_N'(1)G_X'(1) = \E(N)\E(X_1)$.
\end{proof}

\begin{example}[复合泊松：年度保险
损失]\label{ex:b2:genfun:compoundpoisson}
一家保险公司一年内收到 $N \sim \mathcal P(\lambda)$ 次
理赔，每次理赔费用为 $X_k$（以整数为单位，独立同分布，
概率母函数为 $G_X$，均值为 $\mu$，且与 $N$ 独立）。由
\cref{thm:b2:genfun:compound}，总损失 $S$ 满足
\[
G_S(t) = \eu^{\lambda(G_X(t) - 1)},
\qquad
\E(S) = \lambda\mu ,
\]
并在 $1^-$ 处两次求导：
\[
V(S) = \lambda\,G_X''(1) + \lambda^2\mu^2 + \lambda\mu -
(\lambda\mu)^2 = \lambda\,\E(X^2) .
\]
方差涉及单次理赔的\emph{二阶}矩，而不是方差本身：复合泊松和
会对偶发的大额理赔“计入两次”——一次在次数上，
一次在金额上。
取 $\lambda = 10$，理赔金额服从几何分布，
均值为 $2$（$\E X^2 = 6$），则 $\E S = 20$，$V(S) = 60$，
而切比雪夫不等式（\cref{ch:b2:randomvar}）已足以给出可用的
偿付能力边际。“随机停止和”的这一模式，
也将推动
\cref{prop:b2:genfun:branching}中的分支递推：母函数的复合
正是随机种群的代数。
\end{example}

\begin{remark}
$N$ 与各求和项的独立性并非可有可无。
取 $X_k \in \{0, 2\}$ 且两值等概率，令 $N =
X_1$（显然依赖）：则 $S = X_1 + \dots + X_N$ 为
$0$（当 $X_1 = 0$ 时），$2 + X_2$（当 $X_1 = 2$ 时），故 $\E(S) =
\frac12(2 + 1) = \frac32$，而 $\E(N)\E(X_1) = 1\cdot1 =
1$，瓦尔德恒等式不成立。当项数被允许
对各项本身\emph{作出反应}时，干净的乘积结构便
不复存在——此类“停止”规则的完整理论
是本科三年级分册中的鞅论
一章。
\end{remark}

\section{泊松逼近}

\begin{theorem}[稀有事件律]\label{thm:b2:genfun:rareevents}
设 $X_n \sim \mathcal{B}(n, p_n)$，且 $n\,p_n \to \lambda > 0$。
则对每个 $k \in \N$，有
\[
\P(X_n = k)
\xrightarrow[n\to\infty]{}
e^{-\lambda}\frac{\lambda^k}{k!} :
\]
即许多稀有独立事件的二项分布收敛于
参数为 $\lambda$ 的泊松分布。
\end{theorem}

\begin{proof}
直接计算：取 $p_n = \frac{\lambda_n}{n}$，$\lambda_n
\to \lambda$:
\[
\P(X_n = k)
= \binom nk p_n^k(1 - p_n)^{n-k}
= \frac{n(n-1)\cdots(n-k+1)}{n^k}\cdot
\frac{\lambda_n^k}{k!}\,
\bigl(1 - \tfrac{\lambda_n}{n}\bigr)^{n-k} .
\]
当 $n \to \infty$ 而 $k$ 固定时：第一个因子趋于 $1$
（$k$ 个因子之积 $\to 1$）；$\lambda_n^k \to \lambda^k$；
而 $\bigl(1 - \frac{\lambda_n}{n}\bigr)^{n-k} =
\exp\bigl((n-k)\ln(1 - \frac{\lambda_n}{n})\bigr) \to
e^{-\lambda}$，这是因为 $(n - k)\ln\bigl(1 - \frac{\lambda_n}{n}\bigr)
\sim -\lambda_n \to -\lambda$
（见\cref{ch:b2:comparison}）。另一种做法是
在母函数层面：$G_{X_n}(t) = \bigl(1 +
\frac{\lambda_n(t-1)}{n}\bigr)^n \to e^{\lambda(t - 1)} =
G_{\mathcal{P}(\lambda)}(t)$，对每个固定的 $t \in [0, 1]$ ---
即母函数的收敛（对 $\N$ 值随机变量），
而后者等价于每个 $\P(X_n = k)$ 的收敛；
见\cref{exo:b2:genfun:9}。
\end{proof}

\begin{remark}
这正是泊松分布刻画稀有事件计数的原因——例如每页的错字数、
每秒的放射性衰变数、每天在某路口的事故数：
每个机会几乎可以忽略，机会数量却很多，
而极限中只有平均速率 $\lambda$ 保留下来。
\end{remark}

\begin{example}[观察泊松极限
的收敛]\label{ex:b2:genfun:poissonnumerics}
取 $\lambda = 2$，令 $X_n \sim \mathcal B(n, 2/n)$。
无事件概率恰为 $\P(X_n = 0) = (1 - 2/n)^n$：
\[
n = 10:\ 0.107, \qquad
n = 20:\ 0.122, \qquad
n = 50:\ 0.130, \qquad
n = 100:\ 0.133,
\]
与极限 $\eu^{-2} \approx 0.135$ 相比。收敛
是单调的，速度为 $O(1/n)$ --- 展开后 $(1 -
2/n)^n = \eu^{-2}\bigl(1 - \tfrac2n + O(n^{-2})\bigr)$ ---
因此当 $n$ 为数百时，泊松模型已
精确到小数点后第三位。这正是
稀有事件律的实际内涵：建模者从不单独知道 $n$ 与
$p$（一页纸上可能出错的微小机会有多少个？），
只知道它们的乘积 $\lambda$；而极限律幸运地
不依赖于其他任何量。
\end{example}

\section{分支过程}

考虑从一个祖先出发的种群；每个个体
独立地，子代数为随机变量，其分布律为 $(p_k)_{k
\in \N}$，母函数为 $G$（\emph{子代
分布}）。设 $Z_n$ 为第 $n$ 代的大小（$Z_0 =
1$），令 $m = G'(1) = \E(Z_1)$ 为平均子代数。

\begin{proposition}\label{prop:b2:genfun:branching}
$Z_n$ 的母函数是第 $n$ 次迭代
$G_{Z_n} = G \circ G \circ \dots \circ G$（$n$ 次复合），且
灭绝概率 $q_n = \P(Z_n = 0)$ 满足
\[
q_0 = 0, \qquad q_{n+1} = G(q_n),
\]
并单调递增趋于最终灭绝的概率 $q$，它
是 $G$ 的不动点。
\end{proposition}

\begin{proof}
第 $n + 1$ 代是 $Z_n$ 个第 $n$ 代成员的子代数之随机和，
各计数彼此独立，且与 $Z_n$ 独立：
由\cref{thm:b2:genfun:compound}可得 $G_{Z_{n+1}} =
G_{Z_n} \circ G$，并由 $G_{Z_0}(t) = t$ 归纳得到
$n$ 重迭代 --- 由复合运算的结合律，它也
可写作 $G_{Z_{n+1}} = G \circ G_{Z_n}$。对后一形式在
$0$ 处求值：$q_{n+1} = G_{Z_{n+1}}(0) =
G\bigl(G_{Z_n}(0)\bigr) = G(q_n)$。事件 $\{Z_n = 0\}$ 单调递增（已灭绝的种群保持灭绝），
因此 $q_n \uparrow q = \P\bigl(\bigcup_n\{Z_n = 0\}\bigr)$
由单调连续性（见\cref{thm:b2:proba:continuity}），
又 $G$ 在 $[0, 1]$ 上的连续性使 $q_{n+1} = G(q_n)$ 在极限下化为 $q =
G(q)$。
\end{proof}

\begin{example}[观察灭绝概率的收敛]
\label{ex:b2:genfun:cobwebnumerics}
对子代分布 $(p_0, p_1, p_2) = (\tfrac14,
\tfrac14, \tfrac12)$（见\cref{ex:b2:genfun:branchingexample}），
$G(t) = \tfrac14 + \tfrac14t + \tfrac12t^2$，且迭代
$q_{n+1} = G(q_n)$ 给出
\[
q_1 = 0.25, \quad q_2 = 0.34375, \quad q_3 \approx 0.39502,
\quad q_4 \approx 0.42678, \quad q_5 \approx 0.44776,
\]
逐步逼近灭绝概率 $q = \tfrac12$。
差距 $q - q_n$ 依次为 $0.25$、$0.156$、$0.105$、$0.073$、
$0.052$：每一项约为前一项的 $\tfrac34$，且
事实上由中值定理可得 $q - q_{n+1} = G'(c_n)(q
- q_n)$，其中 $G'(q) = \tfrac14 + q = \tfrac34$。由此有两点启示：
一支仍存活于第 $n$ 代的家系，在同一计算中
内含其此后注定灭绝的概率 $q - q_n$；
另一方面，下图中阶梯的收敛速率
恰为不动点处的导数 ---
周末问题将这两点观察化为定理。
\end{example}

\begin{theorem}[灭绝判别准则]\label{thm:b2:genfun:extinction}
设 $p_1 \neq 1$。灭绝概率 $q$ 为 $G$ 在 $\intcc{0}{1}$ 中的
\emph{最小}不动点，且：
\begin{itemize}
\item 若 $m \leq 1$（次临界或临界），则 $q = 1$：灭绝
必然发生；
\item 若 $m > 1$（超临界），则 $q < 1$：种群以正概率
$1 - q$ 永久存续。
\end{itemize}
\end{theorem}

\begin{proof}
$G$ 在 $\intcc{0}{1}$ 上为凸函数（系数非负的幂级数：$G'' \geq 0$），
单调不减，且 $G(1) = 1$。

\emph{最小不动点：}设 $r \in \intcc{0}{1}$ 为任一
不动点。于是 $q_0 = 0 \leq r$，且归纳地 $q_{n+1} = G(q_n)
\leq G(r) = r$（由单调性）；故 $q = \lim q_n \leq r$。

\emph{情形 $m \leq 1$：}假设 $r < 1$ 为一不动点。由
$[r, 1]$ 上的中值定理，存在 $c \in \intoo{r}{1}$，使得
$G'(c) = \frac{G(1) - G(r)}{1 - r} = \frac{1 - r}{1 - r} = 1$。但
$G'$ 单调不减（凸性），且 $\lim_{t\to1^-}G'(t) = m
\leq 1$，故 $G' \leq 1$ 在 $\intoo{0}{1}$ 上；等式 $G'(c) =
1$，于是迫使 $G'$ 恒等于 $1$，从而
在 $\intco{c}{1}$ 上 $G'' = \sum n(n-1)p_nt^{n-2} \equiv 0$。
系数非负的幂级数若在某区间上恒为零，则这些系数全为零：
$p_n = 0$（$n \geq 2$），故 $G(t) = p_0
+ p_1t$，且 $1 = G'(c) = p_1$ --- 与假设
$p_1 \neq 1$ 矛盾。故 $1$ 为唯一不动点，$q = 1$。

\emph{情形 $m > 1$：}在 $1$ 附近，$G(t) - t$ 的导数为 $G'(t) -
1 \to m - 1 > 0$（$t \to 1^-$），于是 $G(t) - t < G(1) - 1 = 0$ 在
某区间 $\intoo{1 - \delta}{1}$ 上：连续函数
$G(t) - t$ 为 $\geq 0$（$t = 0$ 时，$G(0) = p_0 \geq 0$），而 $< 0$
于 $1$ 略小处，由介值定理知它在某 $r < 1$ 处为零。
因此最小不动点满足 $q \leq r < 1$。
\end{proof}

\begin{figure}[ht]
\centering
% Subcritical: G(t) = 0.5 + 0.1 t + 0.4 t^2 => m = 0.9 <= 1
\begin{tikzpicture}[scale=2.9]
  \draw[->] (0,0) -- (1.15,0) node[below] {$t$};
  \draw[->] (0,0) -- (0,1.15);
  \draw[dashed] (0,0) -- (1,1);
  \draw[thick, ocBlue, domain=0:1, smooth, samples=40]
    plot (\x, {0.5 + 0.1*\x + 0.4*\x*\x});
  \fill (1,1) circle (0.02);
  \node[below, font=\small] at (0.55,-0.06)
    {$m \leq 1$：$q = 1$};
  % cobweb from 0
  \draw[ocRed, ->] (0,0) -- (0,0.5);
  \draw[ocRed, ->] (0,0.5) -- (0.5,0.5);
  \draw[ocRed, ->] (0.5,0.5) -- (0.5,0.65);
  \draw[ocRed, ->] (0.5,0.65) -- (0.65,0.65);
  \draw[ocRed, ->] (0.65,0.65) -- (0.65,0.734);
  \draw[ocRed, ->] (0.65,0.734) -- (0.734,0.734);
  \node[font=\small, ocBlue] at (0.32,0.75) {$G$};
\end{tikzpicture}
\qquad
% Supercritical: G(t) = 0.25 + 0.15 t + 0.6 t^2 => m = 1.35 > 1
\begin{tikzpicture}[scale=2.9]
  \draw[->] (0,0) -- (1.15,0) node[below] {$t$};
  \draw[->] (0,0) -- (0,1.15);
  \draw[dashed] (0,0) -- (1,1);
  \draw[thick, ocBlue, domain=0:1, smooth, samples=40]
    plot (\x, {0.25 + 0.15*\x + 0.6*\x*\x});
  \fill (1,1) circle (0.02);
  % fixed point q: 0.6q^2 - 0.85q + 0.25 = 0 -> q = 5/12 ~ 0.4167
  \fill[ocRed] (0.4167,0.4167) circle (0.02)
    node[below right, font=\small] {$q$};
  \node[below, font=\small] at (0.55,-0.06)
    {$m > 1$：$q < 1$};
  % cobweb from 0
  \draw[ocRed, ->] (0,0) -- (0,0.25);
  \draw[ocRed, ->] (0,0.25) -- (0.25,0.25);
  \draw[ocRed, ->] (0.25,0.25) -- (0.25,0.325);
  \draw[ocRed, ->] (0.25,0.325) -- (0.325,0.325);
  \draw[ocRed, ->] (0.325,0.325) -- (0.325,0.3623);
  \draw[ocRed, ->] (0.325,0.3623) -- (0.3623,0.3623);
  \node[font=\small, ocBlue] at (0.28,0.72) {$G$};
\end{tikzpicture}
\caption{灭绝概率作为不动点迭代，
$q_{n+1} = G(q_n)$，自 $q_0 = 0$ 出发（红色阶梯）。左图为
次临界子代分布 --- 凸曲线位于对角线上方，
迭代爬升至唯一不动点 $1$。右图为超临界分布 --- 曲线在
与对角线相交于
$q < 1$ 处，迭代于此停止：存活概率为 $1 -
q > 0$。}
\label{fig:b2:genfun:branching}
\end{figure}

\begin{remark}[如何读懂蛛网图]
图中，竖直移动施用 $G$（由 $(q_n, q_n)$
竖直上升至 $(q_n, G(q_n))$），水平移动至对角线，
把输出转化为输入：阶梯折线\emph{正是}
递推式 $q_{n+1} = G(q_n)$。$G$ 的凸性与 $G(1) = 1$
只留下两种几何情形。一种是曲线在
$\intco01$ 上始终位于对角线上方（均值 $m \leq 1$）：
阶梯折线在 $1$ 之前无处停止。另一种情形是曲线在某处 $q <
1$ 与对角线相交（$m > 1$）：阶梯折线被困于交点之下，
并收敛到该交点，收敛速率为 $G'(q) < 1$
（见\cref{ex:b2:genfun:cobwebnumerics}）。绝灭定理的全部分析
都在这一幅图中一目了然 ---
这正是值得先画图的原因，
再做计算。
\end{remark}

\begin{example}\label{ex:b2:genfun:branchingexample}
后代分布：无子女、一个子女、两个子女，概率分别为
$\frac14, \frac14, \frac12$。于是 $m = \frac14 + 1 =
\frac54 > 1$，且 $G(t) = \frac14 + \frac14 t + \frac12 t^2$。不动点
满足 $\frac12 t^2 - \frac34 t + \frac14 = 0$，即\ $2t^2 - 3t
+ 1 = (2t - 1)(t - 1) = 0$：$q = \frac12$。该家系
以概率 $\frac12$ 绝灭 --- 而以概率 $\frac12$
它永远存续。
\end{example}

\begin{remark}[本卷中的视角]
本章是全书的交汇点，各要素分别来自
具名之处：级数代数取自
\cref{ch:b2:series}与\cref{ch:b2:powerseries}，概率论取自
\cref{ch:b2:proba}（由单调连续性
证明 $q_n \uparrow q$），以及\cref{ch:b2:randomvar}（$G_X =
\E(t^X)$ 为期望，乘性即乘积定理），
凸性则来自\cref{ch:b2:realfun}至
\cref{ch:b2:affine}。甚至重尾的反常情形也彼此相连：
前一章的圣彼得堡随机变量满足
$G(t) = \sum_k2^{-k}t^{2^k}$，是一个完全收敛的级数
在 $\intcc01$ 上，其在 $1^-$ 处的导数发散 ---
期望无穷，一望即知。一个对象，集本年全部工具于一身：
堪称恰当的末章。
\end{remark}

\begin{remark}[常见误区]
(i) 母函数仅适用于取值于 $\N$ 的变量：
对于带符号或非整数的变量，对象 $\E(t^X)$
失去其幂级数结构（本科三年级分册将用
适用于 $\R$ 的变换取代它）。(ii) 对所算的 $G$ 的首要检验是
$G(1) = 1$；第二项检验是系数
须非负 --- 出现负系数说明是代数失误，
而非新的定律。(iii) 在随机和中，
复合顺序很关键：$G_S = G_N \circ G_X$，
\emph{外层}函数负责计数各项，
反过来复合则毫无意义（$G_X \circ G_N$ 将计数
项的项）。(iv) 乘性要求独立，且随机性来源不同：
$G_{2X}(t) = G_X(t^2)$，而非
$G_X(t)^2$。(v) 在 $1$ 处求导是一种边界
运算：当收敛半径恰为 $1$ 时，如
\cref{ex:b2:genfun:heavytail}，$G'(1^-)$ 可能为无穷，
而矩定理的单调极限表述并非
咬文嚼字的细节，而是诚实的陈述。
\end{remark}

\section*{全卷结语}

母函数是本书一个恰当的收尾对象：
它同时是幂级数（\cref{ch:b2:powerseries}），一种
可和族的工具（\cref{ch:b2:series}），一种期望
（\cref{ch:b2:randomvar}），一个凸函数，其几何决定
绝灭（\cref{ch:b2:realfun}），以及不动点迭代
（\cref{ch:b2:metric}）。本科二年级的数学本是一门学科。
本科三年级分册将开启此处有意关闭的
大门：勒贝格积分（补足
\cref{ch:b2:integration}中的控制收敛定理），测度论框架下的
不可数空间概率，以及反函数
定理的完整证明（\cref{ch:b2:diffcalc}），置于
微分几何的框架之中。

\section{练习}

\begin{exercise}[$\star$]\label{exo:b2:genfun:1}
求 $\{1, 2,
\dots, 6\}$ 上均匀分布（一颗均匀骰子）的母函数。证明两颗均匀骰子之和
\emph{不可能}在 $\{2, \dots, 12\}$ 上均匀分布：分解 $G_{X+Y}$
并数根。\emph{（若和均匀，则须有 $G_X(t)G_Y(t) =
\frac{t^2}{11}\sum_{k=0}^{10}t^k$，其非零根为
$11$ 次单位根，除 $1$ 外，无一为实数 ---
而 $G_X/t$ 与 $G_Y/t$ 均为次数为 $5$ 的实系数多项式，
各自至少有一个实根。）}
\end{exercise}

\begin{exercise}[$\star$]\label{exo:b2:genfun:2}
利用母函数，求二项分布与
几何分布的 $\E$ 与 $V$（见\cref{thm:b2:genfun:moments}）。
\end{exercise}

\begin{exercise}[$\star$]\label{exo:b2:genfun:3}
两颗灌铅骰子：能否将两颗骰子（独立地，
同样或不同地）灌铅，使其和在 $\{2, \dots,
12\}$ 上均匀？\emph{（与\cref{exo:b2:genfun:1}中相同的分解障碍：
即使灌铅不同，答案仍是否定的，
因为每个因子 $G_X(t)/t$ 次数为奇数 $5$，
从而有实根，而目标分布无实根。）}
\end{exercise}

\begin{exercise}[$\star\star$]\label{exo:b2:genfun:4}
设 $X_1, X_2, \dots$ 为独立的伯努利变量 $\mathcal{B}(p)$
且 $N \sim \mathcal{P}(\lambda)$ 与它们独立。
利用\cref{thm:b2:genfun:compound}，证明 $S = X_1 + \dots + X_N
\sim \mathcal{P}(\lambda p)$：泊松个数的项，每项
以概率 $p$ 保留，则仍得泊松个数 --- \emph{稀疏化}。
另求被舍弃项数的分布，并为之赞叹：
它是 $\mathcal{P}(\lambda(1-p))$，并可证明它
与 $S$ 独立。
\end{exercise}

\begin{exercise}[$\star\star$]\label{exo:b2:genfun:5}
（负二项分布）设 $T_r$ 为获得
$r$ 次正面（正面概率为 $p$）所需的抛掷次数。将 $T_r$ 写成 $r$ 个
独立几何变量之和，并由此推出
\[
G_{T_r}(t) = \Bigl(\frac{pt}{1 - (1-p)t}\Bigr)^{r},
\qquad
\E(T_r) = \frac rp,
\qquad
V(T_r) = \frac{r(1-p)}{p^2},
\]
并展开 $G_{T_r}$ 以求 $\P(T_r = n) = \binom{n-1}{r-1}
p^r(1-p)^{n-r}$。
\end{exercise}

\begin{exercise}[$\star\star$]\label{exo:b2:genfun:6}
对后代分布 $p_0 = \frac18$，$p_1 = \frac38$，$p_2 =
\frac38$，$p_3 = \frac18$：求 $m$，判断
是否超临界，并求灭绝概率 $q$
精确地求出。\emph{（分解出 $t = 1$ 这一根，即 $G(t) - t$ 的因式。）}
\end{exercise}

\begin{exercise}[$\star\star\star$]\label{exo:b2:genfun:7}
（总后代）在次临界分支过程（$m < 1$）中，设
$Y = \sum_{n\geq0} Z_n$ 为曾经出生的
个体总数。证明 $\E(Y) = \sum_n m^n = \frac{1}{1 - m}$（论证
求和交换的合理性），并证明概率母函数 $H = G_Y$ 满足函数方程
$H(t) = t\,G(H(t))$。
\emph{（祖先，加上其每个子代的
总后代，它们是 $Y$ 的独立副本。）}
\end{exercise}

\begin{exercise}[$\star\star\star$]\label{exo:b2:genfun:8}
设 $X$ 的概率母函数为 $G$，其收敛半径 $>
1$。证明\emph{指数尾界}：存在 $C > 0$ 与
$\rho \in \intoo{0}{1}$，使得 $\P(X \geq n) \leq C\rho^n$。
\emph{（将马尔可夫
不等式用于 $t^X$，$t > 1$ 为收敛圆内某固定值。）}反之，证明若 $\P(X \geq n) \leq C\rho^n$ 且
$\rho < 1$，则 $G$ 的收敛半径 $\geq 1/\rho > 1$。
\end{exercise}

\begin{exercise}[$\star\star\star$]\label{exo:b2:genfun:9}
（连续性定理，初等情形）设 $X, X_1, X_2, \dots$ 为
$\N$ 值的随机变量，且 $G_{X_n}(t) \to G_X(t)$ 对每个 $t \in
\intco{0}{1}$ 成立。证明 $\P(X_n = k) \to \P(X = k)$ 对每个
$k$ 成立。\emph{（对 $k$ 归纳：当 $k = 0$ 时取 $t \to 0$ ---
注意：固定较小的 $t$，利用 $\abs{\P(X_n = 0) - G_{X_n}(t)}
\leq \frac{t}{1-t}$，这是因为尾部 $\sum_{j \geq 1}p_jt^j
\leq \frac{t}{1 - t}$ 成立；再用对角线法。归纳步骤中，
考虑 $\frac{G(t) - \P(X = 0)}{t}$，它是平移后
分布的概率母函数。）}
\end{exercise}

\begin{exercise}[$\star$]\label{exo:b2:genfun:10}
（奇偶技巧）对取值于 $\N$ 的随机变量 $X$，证明
\[
\P(X \text{ 为偶数}) = \frac{1 + G_X(-1)}{2} ,
\]
并对 $X \sim \mathcal P(\lambda)$ 与
$X \sim \mathcal B(n, p)$ 两种情形计算此概率。$G_X(-1) \to 0$ 从概率上
意味着什么？
\end{exercise}

\begin{exercise}[$\star\star$]\label{exo:b2:genfun:11}
（西彻曼骰子）验证公平骰子的
概率母函数的分解式为
\[
\frac{t + t^2 + \dots + t^6}{6}
= \frac{t\,(1 + t)(1 + t + t^2)(1 - t + t^2)}{6},
\]
并证明面为 $\{1, 2, 2, 3, 3, 4\}$
与 $\{1, 3, 4, 5, 6, 8\}$ 的两颗骰子，其概率母函数分别为
$\frac{t(1+t)(1+t+t^2)}6$ 和
$\frac{t(1+t)(1+t+t^2)(1-t+t^2)^2}6$，二者之积
与两颗标准骰子之积相同：这些奇异骰子产生每个点数和 $2,
\dots, 12$ 的概率与标准骰子完全相同。
\end{exercise}

\begin{exercise}[$\star\star\star$]\label{exo:b2:genfun:12}
（等待连续两次正面）一枚正面概率为
$p$ 的硬币反复抛掷，直到连续出现两次正面为止；设 $T$ 为
抛掷次数（即\cref{exo:b2:proba:6}中的游戏）。
以前几次抛掷的结果为条件，建立
由状态“当前无正面”与“当前有一个正面”出发的母函数所满足的线性方程组，
并得出
\[
G_T(t) = \frac{p^2t^2}{1 - qt - pqt^2}
\qquad (q = 1 - p);
\]
验证 $G_T(1) = 1$ 及 $\E(T) = \dfrac{1 + p}{p^2}$（对公平硬币
$= 6$）。
\end{exercise}

\section{问题：高尔顿–沃森过程的求解}

\begin{problem}[{周末问题 --- 增长速率、精确
解、总后代与柯尔莫哥洛夫的临界估计
}]\label{pb:b2:genfun:1}
灭绝准则（见\cref{thm:b2:genfun:extinction}）
将分支过程\index{分支过程}划分为
次临界、临界与超临界 --- 但它
没有涉及\emph{速率}：注定灭绝的分支以多快的速度消亡，
幸存的分支又能增长到多大。本题计算这些速率。沿用本章记号：
后代分布 $(p_k)$，其概率母函数为 $G$，
均值 $m = G'(1)$，各代规模 $Z_n$（$Z_0 = 1$），迭代
$G_n = G_{Z_n}$，灭绝概率 $q_n = \P(Z_n = 0)
\uparrow q$；始终假定 $p_1 \neq 1$；若出现二阶矩，
则假定 $G''(1) < \infty$，并记 $\sigma^2 =
V(Z_1)$。

\medskip
\noindent\textbf{第 I 部分 --- 各代的矩。}
\begin{enumerate}
  \item 证明 $\E(Z_n) = m^n$ \emph{（对 $G_n = G
        \circ G_{n-1}$ 在 $1^-$ 处应用链式法则，利用 $G_{n-1}(1) = 1$ 及
        \cref{thm:b2:genfun:moments}）}。
  \item 建立递推式 $G_n''(1) =
        G''(1)\,m^{2(n-1)} + m\,G_{n-1}''(1)$，并解之：
        $G_n''(1) = G''(1)\,m^{n-1}\dfrac{m^n - 1}{m - 1}$
        当 $m \neq 1$ 时为上式，又 $G_n''(1) = n\,G''(1)$ 对 $m =
        1$ 时成立。
  \item 推出
        \[
        V(Z_n) = \sigma^2m^{n-1}\,\frac{m^n - 1}{m - 1}
        \quad (m \neq 1),
        \qquad
        V(Z_n) = n\,\sigma^2 \quad (m = 1).
        \]
  \item （次临界速率，上界）对 $m < 1$，证明
        $\P(Z_n > 0) \leq m^n$ \emph{（对取整数值的
        $Z_n$ 用马尔可夫不等式）}：灭绝必然发生，
        且以几何速率进行 --- 是对本章判别准则的
        定量加强。
  \item （次临界速率，下界）利用
        柯西--施瓦茨不等式作用于 $Z_n\mathbf 1_{Z_n > 0}$，证明
        \[
        \P(Z_n > 0) \geq \frac{\E(Z_n)^2}{\E(Z_n^2)}
        \geq c\,m^{n}
        \quad\text{其中}\quad
        c = \Bigl(\frac{\sigma^2}{m(1-m)} + 1\Bigr)^{-1} :
        \]
        几何速率 $m^n$ 精确到常数因子为止。
\end{enumerate}

\medskip
\noindent\textbf{第 II 部分 --- 几何族，
精确求解。} 设后代分布为 $\N$ 上的几何分布：$p_k =
qp^k$（$k \geq 0$），其中 $0 < p < 1$，$q = 1 - p$。
\begin{enumerate}[resume]
  \item 计算 $G(t) = \dfrac{q}{1 - pt}$ 与 $m = \dfrac
        pq$，并用 $p$ 刻画三种情形。
  \item 解方程 $G(t) = t$，证明不动点为 $1$ 与
        $q/p = 1/m$，并求得灭绝概率
        $q_{\mathrm{ext}} = \min(1, 1/m)$。
  \item 用归纳法证明下列闭式：
        \[
        q_n = \frac{m^n - 1}{m^{n+1} - 1} \quad (m \neq 1),
        \qquad
        q_n = \frac{n}{n+1} \quad (m = 1).
        \]
  \item 由此得精确速率：次临界情形下 $1 - q_n \sim (1 - m)\,m^n$，
        超临界情形下 $q_{\mathrm{ext}} - q_n
        \sim \dfrac{m - 1}{m^{2}}\cdot m^{-n}$，
        并验证超临界情形的收缩比
        为 $G'(q_{\mathrm{ext}}) = 1/m$。
  \item 临界情形（$p = \tfrac12$）：计算 $\sigma^2 =
        2$，并注意到 $1 - q_n = \frac1{n+1}$：生存概率
        按 $\frac1n$ 的量级衰减 --- 既不是几何型，也不可和。
  \item 仍为临界情形：用归纳法证明完整的迭代式
        \[
        G_n(t) = \frac{n - (n-1)t}{n + 1 - nt},
        \]
        由此推出：在生存条件下，$Z_n$ 为
        $\N^*$ 上参数为 $\frac1{n+1}$ 的几何分布：
        \[
        \P(Z_n = k \mid Z_n > 0) = \frac1{n+1}
        \Bigl(\frac{n}{n+1}\Bigr)^{k-1},
        \qquad
        \E(Z_n \mid Z_n > 0) = n + 1 .
        \]
        平均谱系趋于消亡，但存活的谱系
        规模为 $n$ 量级。
\end{enumerate}

\medskip
\noindent\textbf{第 III 部分 --- 总后代数。}设 $Y =
\sum_{n\geq0}Z_n \in \N^* \cup \{\infty\}$ 为历史上出生的
个体总数，并记 $H(t) =
\sum_{k\geq1}\P(Y = k)t^k$。
\begin{enumerate}[resume]
  \item 证明 $\P(Y < \infty) = q_{\mathrm{ext}}$，并
        回顾\cref{exo:b2:genfun:7}中的函数方程
        $H(t) = t\,G(H(t))$（其推导未用到
        $m < 1$）。
  \item （二元分支）对 $p_0 = p_2 = \frac12$
        （临界情形），解函数方程：
        \[
        H(t) = \frac{1 - \sqrt{1 - t^2}}{t},
        \]
        并利用\cref{ex:b2:powerseries:catalan}展开，
        得到
        \[
        \P(Y = 2k + 1) = \frac{C_k}{2^{2k+1}},
        \qquad C_k = \frac1{k+1}\binom{2k}k ;
        \]
        验证 $\P(Y = 1) = \frac12$ 与 $\P(Y = 3)
        = \frac18$，两者均可直接计数检验。
  \item 在 $1^-$ 处对函数方程求导，
        证明 $\E(Y) = \frac{1}{1-m}$（$m < 1$），而
        临界性则迫使 $\E(Y) = \infty$：临界总后代
        几乎必然有限，但具有无穷
        均值。
  \item 利用中心二项式的渐近式
        （\cref{ex:b2:comparison:centralbinomial}），证明
        \[
        \P(Y = 2k+1) \sim \frac{1}{2\sqrt\pi\,k^{3/2}},
        \]
        即 $k^{-3/2}$ 型的重尾，并推出 $\P(Y > n)
        \asymp n^{-1/2}$（只需该阶的上、下界
        即可）。
  \item 与公平随机游走作比较（周末问题，见
        \cref{ch:b2:proba}）：那里返回必然发生，但
        返回时间均值无穷，这里的总后代数几乎必然有限，
        但均值无穷，两者都带有
        $n^{-3/2}$ 型的局部律。请用一段话说明
        临界性为何产生这一特征。
\end{enumerate}

\medskip
\noindent\textbf{第 IV 部分 --- 临界情形下的柯尔莫哥洛夫
估计。}设 $m = 1$，$0 < \sigma^2 = G''(1) <
\infty$。
\begin{enumerate}[resume]
  \item 证明 $G''$ 可连续延拓至 $\intcc01$
        \emph{（非负递增，且极限有限）}，并
        由此推出 $1$ 处的泰勒展开：
        \[
        G(t) = t + b\,(1-t)^2 + o\bigl((1-t)^2\bigr),
        \qquad b = \frac{G''(1)}2 = \frac{\sigma^2}2 .
        \]
  \item 对 $t \in \intco01$，记 $h(t) = \dfrac1{1 - G(t)} -
        \dfrac1{1 - t}$。证明
        \[
        h(t) = \frac{G(t) - t}{(1 - G(t))(1 - t)}
        \xrightarrow[t\to1^-]{} b .
        \]
  \item 沿迭代序列 $q_{j+1} = G(q_j)$ 裂项求和：
        \[
        \frac1{1 - q_n} = 1 + \sum_{j=0}^{n-1}h(q_j),
        \]
        并用切萨罗平均法得出
        \[
        \P(Z_n > 0) = 1 - q_n \sim \frac{2}{\sigma^2\,n}
        \]
        --- \emph{柯尔莫哥洛夫估计}：每个临界分支过程
        都以普适速率 $1/n$ 消亡，
        仅常数与后代分布有关。
  \item 用第 10 题中的临界几何情形
        检验该估计。
  \item 推出 $\E(Z_n \mid Z_n > 0) = \dfrac{1}{1 - q_n}
        \sim \dfrac{\sigma^2 n}{2}$ \emph{（注意 $\E(Z_n
        \mathbf 1_{Z_n>0}) = \E(Z_n) = 1$）}，并
        与第 11 题核对：在存活条件下，
        种群以\emph{线性}方式增长——
        正走在死亡与爆炸之间的临界钢丝上。
\end{enumerate}

\medskip
\noindent\textbf{第 V 部分 --- 应用与综合。}
\begin{enumerate}[resume]
  \item （传染病、链式反应）对泊松后代分布
        $\mathcal P(\lambda)$ --- 每个病例感染
        $\mathcal P(\lambda)$ 个新病例 --- 写出
        灭绝方程 $q = \eu^{\lambda(q-1)}$，并
        数值求解 $\lambda = 1.5$（$q \approx
        0.417$）和 $\lambda = 2$（$q \approx 0.203$）：
        从一个病例出发，大规模爆发
        即使 $\lambda > 1$ 时也\emph{并非}必然发生。说明
        为何迭代 $q_{n+1} = \eu^{\lambda(q_n - 1)}$
        从 $q_0 = 0$ 出发会收敛到正确的根。
  \item 若从 $k$ 个祖先而非一个出发，证明
        灭绝概率为 $q^k$。应用：
        取 $\lambda = 1.5$，问至少需要多少个初始病例，
        才能使爆发的概率至少为 $99\%$？
  \item （以灭绝为条件的超临界过程）
        设 $m > 1$，灭绝概率为 $q \in
        \intoo01$。先由凸性证明在最小不动点处 $G'(q) < 1$，
        并推出
        $q_{\mathrm{ext}} - q_n = O\bigl(G'(q)^n\bigr)$
        （几何收敛，如第 9 题所示）。
        再证明 $\widehat G(t) = G(qt)/q$ 为某个合法后代分布的概率母函数，
        其均值为 $\widehat m =
        G'(q) < 1$，是一个次临界的伴随过程。
        在几何分布族上验证：以灭绝为条件的超临界
        $(p, q)$ 过程会交换
        $p$ 与 $q$。（完整陈述——条件化后的过程
        \emph{即是}伴随过程——该结论在
        本科三年级分册中证明；这里你已验证了其
        母函数“影子”。）
  \item 综合：列出三分表——对 $m <
        1$、$m = 1$、$m > 1$ 时，$q$ 的取值；$\P(Z_n >
        0)$ 或 $q - q_n$ 的收敛速率；$\E(Y)$；存活
        世代的规模。用一句话分别说明各工具如何
        母函数的复合、凸性、在 $1^-$ 处的泰勒展开与
        切萨罗平均如何贯穿整个问题，以及
        本科三年级分册补充了哪些内容（鞅 $Z_n/m^n$ 与
        亚格洛姆的指数极限律）。
\end{enumerate}
\end{problem}
